\begin{abstract}
Whether a model looks at the image is usually inferred from the average accuracy drop when the image is removed. That average hides when nonvisual evidence becomes so predictive that additional visual evidence stops helping. We study this competition in multiple-choice VQA, where a text-side shortcut can settle answers without the image. We define a benchmark's \emph{visual headroom} as the accuracy a model gains from the true image over its \emph{blind ceiling}, the best accuracy reached without the image or from option text alone. We measure headroom with counterfactual images, a ladder of image strength, and a length shortcut planted at graded strength, so that shortcut strength and image strength vary jointly. Planting the shortcut lowers the image benefit on PMC-VQA from $+13.5$ to $+0.1$ points. The strength at which the image stops mattering depends on the benchmark: $0.85$ for Qwen-3B on PMC-VQA and $0.99$ on A-OKVQA. A boundary prediction fixed before its runs held for two further models, and three attempts to explain the difference between benchmarks failed. Without planting, headroom is $-0.1$ points on Med-CMR and $+4.9$ on OmniMedVQA, against $+13.7$ and $+33.2$ on PMC-VQA and A-OKVQA. On Med-CMR the shortest option is correct for $91.9\%$ of questions. For Qwen2.5-VL-3B, fine-tuning on the benchmark's own training split installed this shortcut within about 1{,}000 questions (two seeds, one benchmark) and left no measurable headroom. A penalty on confident answers under a mismatched image restores measurable image dependence for $4$ to $6$ points of accuracy. The restored dependence is partly coarse: the same penalty loses $11$ to $18$ points against a plausible wrong image. We report failed and post-hoc tests, and all tests follow dated, hash-chained amendments.
\end{abstract}
